% Lista de exercícios — gera lista.pdf e, com -jobname=lista-solucao,
% lista-solucao.pdf (ver tex/exercicios.sty).
\documentclass[lista]{tmlnotas}
\usepackage{exercicios}

\begin{document}
\aula{2}{Computação aplicada: de Python ao aprendizado de máquina}[Computação aplicada]

\noindent Resolva os exercícios em um notebook no Colab. Os Exercícios 2 a 6
usam a base do experimento NSW, lida como no notebook da aula:
\begin{codigo}
import numpy as np
import pandas as pd
URL = "https://users.nber.org/~rdehejia/data/"
nsw = pd.read_stata(URL + "nsw_dw.dta")
cps = pd.read_stata(URL + "cps_controls.dta")
COVARIAVEIS = ["age", "education", "black", "hispanic", "married",
               "nodegree", "re74", "re75"]
\end{codigo}
Níveis: \nivel{1}~básico, \nivel{2}~intermediário, \nivel{3}~desafio.
\ifsolucao\textbf{Esta versão contém as soluções.}\fi

% ---------------------------------------------------------------------
\begin{exercicio}[Fichas de participantes]{1}
Use a lista \py{pessoas} do notebook (cinco dicionários com as chaves
\py{"idade"}, \py{"educ"}, \py{"tratado"} e \py{"re78"}).
\begin{enumerate}[label=(\alph*)]
  \item Acrescente um sexto participante: controle, 19 anos, 8 anos de estudo e
        renda zero em 1978.
  \item Escreva uma função \py{resumo(pessoas, tratado)} que devolva um
        \emph{dicionário} com o número de pessoas do grupo, a renda média e a
        fração sem renda em 1978. O argumento \py{tratado} é \py{True} ou
        \py{False}.
  \item Use a função para os dois grupos e imprima os resultados com
        \emph{f-strings}, a renda com duas casas decimais.
\end{enumerate}
\end{exercicio}
\begin{solucao}
\begin{codigo}
pessoas = [
    {"idade": 37, "educ": 11, "tratado": True,  "re78": 9930.05},
    {"idade": 22, "educ": 9,  "tratado": True,  "re78": 3595.89},
    {"idade": 30, "educ": 12, "tratado": True,  "re78": 24909.45},
    {"idade": 34, "educ": 12, "tratado": False, "re78": 2113.72},
    {"idade": 24, "educ": 10, "tratado": False, "re78": 7618.64},
]
pessoas.append({"idade": 19, "educ": 8, "tratado": False, "re78": 0.0})

def resumo(pessoas, tratado):
    rendas = [p["re78"] for p in pessoas if p["tratado"] == tratado]
    return {"n": len(rendas),
            "renda_media": sum(rendas) / len(rendas),
            "sem_renda": sum(r == 0 for r in rendas) / len(rendas)}

for grupo in [True, False]:
    r = resumo(pessoas, grupo)
    print(f"tratado={grupo}: n={r['n']}, média={r['renda_media']:,.2f}, "
          f"sem renda={r['sem_renda']:.0%}")
\end{codigo}
Os controles passam a ter 3 pessoas, renda média de US\$ 3.244,12 e 33\% sem
renda. Em \py{sum(r == 0 for r in rendas)}, cada comparação vale \py{True} ou
\py{False}, que o Python soma como 1 ou 0.
\end{solucao}

% ---------------------------------------------------------------------
\begin{exercicio}[Filtros e agrupamentos]{1}
\begin{enumerate}[label=(\alph*)]
  \item Com \py{groupby}, calcule a fração de tratados e de controles que
        tiveram renda zero em 1978. O programa reduziu o desemprego?
  \item Restrinja a base a quem tinha renda zero em 1975. Quantas pessoas são?
        Qual a diferença de médias da renda de 1978 entre tratados e controles
        nesse subgrupo?
\end{enumerate}
\end{exercicio}
\begin{solucao}
\begin{codigo}
(nsw["re78"] == 0).groupby(nsw["treat"]).mean()   # 0: 0,354 | 1: 0,243

sub = nsw[nsw["re75"] == 0]
medias = sub.groupby("treat")["re78"].mean()
len(sub), medias[1] - medias[0]                    # 289 pessoas, ~ 1.711
\end{codigo}
(a) Sem renda em 1978: 35,4\% dos controles e 24,3\% dos tratados, uma queda
de cerca de 11 pontos percentuais. Como o tratamento foi sorteado, a diferença
pode ser lida como efeito do programa sobre ter alguma renda. (b) São 289
pessoas; a diferença é de cerca de US\$ 1.711. Como o filtro usa uma variável
anterior ao programa, o subgrupo continua experimental.
\end{solucao}

% ---------------------------------------------------------------------
\begin{exercicio}[Balanceamento padronizado]{2}
A diferença de médias de uma covariável depende da sua unidade (anos, dólares).
Uma medida sem unidade é a \emph{diferença padronizada}
\[
  d_k = \frac{\bar x_{k,1} - \bar x_{k,0}}{\sqrt{(s_{k,1}^2 + s_{k,0}^2)/2}},
\]
em que $\bar x_{k,g}$ e $s_{k,g}^2$ são a média e a variância da covariável $k$
no grupo $g$.
\begin{enumerate}[label=(\alph*)]
  \item Escreva uma função \py{balanco(df)} que devolva $d_k$ para todas as
        \py{COVARIAVEIS}, usando \py{groupby} ou máscaras.
  \item Aplique-a ao experimento e à base \py{pd.concat([tratados, cps])}.
        Uma regra prática considera preocupante $\abs{d_k} > 0{,}25$. Quais
        covariáveis passam desse limite em cada caso?
\end{enumerate}
\end{exercicio}
\begin{solucao}
\begin{codigo}
def balanco(df):
    g1 = df.loc[df["treat"] == 1, COVARIAVEIS]
    g0 = df.loc[df["treat"] == 0, COVARIAVEIS]
    return (g1.mean() - g0.mean()) / np.sqrt((g1.var() + g0.var()) / 2)

tratados = nsw[nsw["treat"] == 1]
nao_exp = pd.concat([tratados, cps])
pd.DataFrame({"experimento": balanco(nsw),
              "CPS": balanco(nao_exp)}).round(2)
\end{codigo}
No experimento, todas as diferenças são pequenas; a maior em módulo é a de
\texttt{nodegree} ($-0{,}30$), um pouco acima do limite, o que acontece por acaso
em amostras pequenas. Na comparação com a CPS, quase tudo passa do limite:
\texttt{black} (2,43), \texttt{re75} ($-1{,}75$), \texttt{re74} ($-1{,}57$),
\texttt{married} ($-1{,}23$), \texttt{nodegree} (0,90), \texttt{age}
($-0{,}80$) e \texttt{education} ($-0{,}68$). Só \texttt{hispanic} fica abaixo.
\end{solucao}

% ---------------------------------------------------------------------
\begin{exercicio}[O efeito é estatisticamente significante?]{2}
\begin{enumerate}[label=(\alph*)]
  \item Mostre que, para uma \emph{dummy} $D$ com proporção $p$ de uns,
        $\Var(D) = p(1-p)$ e $\Cov(y, D) = p(1-p)(\bar y_1 - \bar y_0)$
        (use as versões amostrais com divisor $n$). Conclua que o coeficiente
        do MQO, $\Cov(y,D)/\Var(D)$, é a diferença de médias.
  \item Calcule o erro-padrão da diferença de médias,
        $\sqrt{s_1^2/n_1 + s_0^2/n_0}$, e a estatística $t$. O efeito é
        significante a 5\%?
\end{enumerate}
\end{exercicio}
\begin{solucao}
(a) Como $D_i^2 = D_i$, $\frac1n\sum D_i^2 - p^2 = p - p^2 = p(1-p)$. Para a
covariância, $\frac1n\sum D_i y_i - p\,\bar y = p\,\bar y_1 - p\,\bar y$, e
$\bar y = p\,\bar y_1 + (1-p)\,\bar y_0$; substituindo,
$p\,\bar y_1 - p^2\bar y_1 - p(1-p)\bar y_0 = p(1-p)(\bar y_1 - \bar y_0)$.
A razão é $\bar y_1 - \bar y_0$.
\begin{codigo}
y1 = nsw.loc[nsw["treat"] == 1, "re78"]
y0 = nsw.loc[nsw["treat"] == 0, "re78"]
dif = y1.mean() - y0.mean()
ep = np.sqrt(y1.var() / len(y1) + y0.var() / len(y0))
dif, ep, dif / ep              # ~ 1794, ~ 671, t ~ 2,67
\end{codigo}
(b) O erro-padrão é de cerca de US\$ 671 e $t \approx 2{,}67 > 1{,}96$: o efeito
é significante a 5\%. O intervalo de 95\% vai de cerca de US\$ 479 a US\$ 3.109.
É um intervalo largo, típico de experimentos com poucas centenas de pessoas.
\end{solucao}

% ---------------------------------------------------------------------
\begin{exercicio}[Gradiente descendente com uma covariável]{2}
Estenda o modelo do notebook para
$\widehat{re78} = a + b\cdot treat + c\cdot education$, com a renda em mil US\$.
\begin{enumerate}[label=(\alph*)]
  \item Escreva o EQM como função de \py{params = jnp.array([a, b, c])}, usando
        uma matriz $\mX$ com colunas de uns, \texttt{treat} e \texttt{education}.
        Obtenha o gradiente com \py{jax.grad} e compare a solução com
        \py{np.linalg.lstsq}.
  \item Rode 500 passos com $\lr = 0{,}01$ e com $\lr = 0{,}005$. O que acontece?
  \item Troque \texttt{education} pela versão \emph{padronizada} (subtraia a
        média e divida pelo desvio-padrão) e rode 100 passos com $\lr = 0{,}5$.
        Compare com o item (b).
\end{enumerate}
\end{exercicio}
\begin{solucao}
\begin{codigo}
import jax, jax.numpy as jnp
y = jnp.array(nsw["re78"].to_numpy() / 1000)

def treinar(X, eta, passos):
    X = jnp.array(X)
    eqm = lambda p: jnp.mean((y - X @ p) ** 2)
    grad = jax.grad(eqm)
    p = jnp.zeros(X.shape[1])
    for _ in range(passos):
        p = p - eta * grad(p)
    return p

um = np.ones(len(nsw))
X = np.column_stack([um, nsw["treat"], nsw["education"]])
np.linalg.lstsq(X, np.asarray(y), rcond=None)[0]   # b ~ 1,692
treinar(X, 0.01, 500)      # diverge (valores enormes)
treinar(X, 0.005, 500)     # b ~ 1,20: ainda longe

z = (nsw["education"] - nsw["education"].mean()) / nsw["education"].std()
Xz = np.column_stack([um, nsw["treat"], z])
treinar(Xz, 0.5, 100)      # b ~ 1,692: convergiu
\end{codigo}
(b) Com $\lr = 0{,}01$ o treino diverge; com $\lr = 0{,}005$ ele não diverge,
mas depois de 500 passos o coeficiente do tratamento ainda está em cerca de 1,20
mil, longe do 1,69 do MQO. (c) Padronizada, a escolaridade fica na mesma escala
das outras colunas, um passo grande passa a ser seguro e 100 passos bastam.
O coeficiente do tratamento é o mesmo, porque padronizar uma covariável não muda
o coeficiente das demais. A Aula~4 explica a geometria por trás disso.
\end{solucao}

% ---------------------------------------------------------------------
\begin{exercicio}[Comparar apenas quem é comparável]{3}
O logit do notebook mostrou que quase todos os homens da CPS têm probabilidade
prevista de ser tratado perto de zero. Uma ideia é descartá-los e comparar os
tratados apenas com quem se parece com eles.
\begin{enumerate}[label=(\alph*)]
  \item Na base \py{nao_exp} (tratados mais CPS), ajuste
        \py{LogisticRegression(max_iter=5000)} de \texttt{treat} nas
        \py{COVARIAVEIS}, com as rendas em mil US\$, na amostra inteira.
        Guarde a probabilidade prevista numa coluna \texttt{p}.
  \item Mantenha só as observações com $p > 0{,}1$. Quantos tratados e quantos
        homens da CPS sobram? Calcule a diferença de médias da renda de 1978 e o
        coeficiente de \texttt{treat} num MQO com as covariáveis.
  \item Repita com $p > 0{,}05$. Compare com o experimento (US\$ 1.794) e
        discuta: o que se ganhou e o que se perdeu?
\end{enumerate}
\end{exercicio}
\begin{solucao}
\begin{codigo}
from sklearn.linear_model import LogisticRegression

X = nao_exp[COVARIAVEIS].copy()
X[["re74", "re75"]] /= 1000
logit = LogisticRegression(max_iter=5000).fit(X, nao_exp["treat"])
nao_exp = nao_exp.assign(p=logit.predict_proba(X)[:, 1])

def avaliar(corte):
    sub = nao_exp[nao_exp["p"] > corte]
    medias = sub.groupby("treat")["re78"].mean()
    Xs = np.column_stack([np.ones(len(sub)), sub["treat"], sub[COVARIAVEIS]])
    b = np.linalg.lstsq(Xs, sub["re78"].to_numpy(), rcond=None)[0][1]
    return sub["treat"].value_counts().to_dict(), medias[1] - medias[0], b

avaliar(0.10)    # 145 tratados, 358 CPS | dif. ~ 1.444 | MQO ~ 1.645
avaliar(0.05)    # 155 tratados, 520 CPS | dif. ~ 669   | MQO ~ 1.600
\end{codigo}
Ao descartar os homens da CPS sem semelhantes entre os tratados, as estimativas
se aproximam muito do experimento: o MQO com covariáveis vai de US\$ 699 (CPS
inteira) para cerca de US\$ 1.600 a 1.645. Mas há custos. (i) Também saem
tratados (sobram 145 ou 155 dos 185), então estimamos o efeito para um
subgrupo. (ii) O resultado depende do corte escolhido: a diferença de médias
simples é de US\$ 1.444 com o corte de 0,1 e de US\$ 669 com o de 0,05.
(iii) Aqui sabemos a resposta certa; numa avaliação real, não saberíamos qual
corte ``acerta''. Métodos que tratam essa escolha com mais cuidado (pareamento,
ponderação e o DML) são assunto da Aula~6.
\end{solucao}

\end{document}
